% Auto-converted from Optimizing_Gestural_Efficiency.docx via pandoc
% Source section: Conclusion
\chapter{Conclusion}\label{conclusion}

\section{Future Work}\label{future-work}
\section{Optimizations: If only one lex per selected prefix,
insert!}\label{optimizations-if-only-one-lex-per-selected-prefix-insert}

\section{Exhaustive Mapping of Unihan Chars to Pinyins using
Unihan}\label{exhaustive-mapping-of-unihan-chars-to-pinyins-using-unihan}

\section{Arranging Characters by Frequency (Might Be Already
Done)}\label{arranging-characters-by-frequency-might-be-already-done}

\section{Reduce Number of Gestures By Phonemes vs. Graphemes To
Lexons}\label{reduce-number-of-gestures-by-phonemes-vs.-graphemes-to-lexons}

Use pīnyīn packed by phonemic analysis to distribute by syllable instead
of pinyin

\section{Relate Frequency List Size to Percentage of Text
Covered}\label{relate-frequency-list-size-to-percentage-of-text-covered}

LCMC -- 5000 entries 66\% - How many gestures to implement?

Other -- 50,000 entries 75\% - How many gestures to implement?

\section{Can GE be increased using different corpora for different
languages?}\label{can-ge-be-increased-using-different-corpora-for-different-languages}

\section{Can overall effort be reduced by rearranging markmaps,
chordmaps?}\label{can-overall-effort-be-reduced-by-rearranging-markmaps-chordmaps}

The current research suggests that although complex gestures incur a
`cost' in time and/or effort, they do not constitute what is intuitively
a `cognitive cost'. Possibly further research can make a testable
hypothesis from this intuition.

\section{Other than rectangular cell arrangement that reduces
eye-scan
time/effort?}\label{other-than-rectangular-cell-arrangement-that-reduces-eye-scan-timeeffort}

\section{Alternative Modes to Minimize time spent for Long
Press}\label{alternative-modes-to-minimize-time-spent-for-long-press}

\section{`Slowly' Rearranging to Catch Eye-Scan
Sooner}\label{slowly-rearranging-to-catch-eye-scan-sooner}

Mode in which when display not visible:

\begin{itemize}
\item
  Short press selects key mark
\item
  Long press selects from initial choice set
\item
  Shift long press displays secondary choice sets with short press
  select and longpress tertiary dispolay
\end{itemize}

Reference

Allen, J. D., Anderson, D., Becker, J., Cook, R., Davis, M., Edberg, P.,
Everson, M., Freytag, A., Iancu, L., \& Ishida, R. (2012a). The unicode
standard. \emph{Mountain view, CA}.

Allen, J. D., Anderson, D., Becker, J., Cook, R., Davis, M., Edberg, P.,
Everson, M., Freytag, A., Iancu, L., \& Ishida, R. (2012b). The unicode
standard. \emph{Mountain view, CA}, 660-664.

Baroni, M., \& Kilgarriff, A. (2006). Large linguistically-processed web
corpora for multiple languages. EACL'06: Proceedings of the Eleventh
Conference of the European Chapter of the Association for Computational
Linguistics: Posters \& Demonstrations; 2006 Apr 5-6; Trento, Italy.
Stroudsburg (PA): Association for Computational Linguistics; 2006. p.
87-90,

Da, J. (2004). A corpus-based study of character and bigram frequencies
in Chinese e-texts and its implications for Chinese language
instruction. Proceedings of the fourth international conference on new
technologies in teaching and learning Chinese,

Francis, G., \& Oxtoby, C. (2006). Building and testing optimized
keyboards for specific text entry. \emph{Hum Factors}, \emph{48}(2),
279-287. \url{https://doi.org/10.1518/001872006777724390}

Hendy, J. (2009). \emph{Graphically enhanced keyboard accelerators for
GUIs} University of British Columbia{]}.

Her, O.-S. (2005). Between globalization and indigenization: On Taiwan's
pinyin issue from the perspectives of the new economy. \emph{Journal of
Humanities and Social Sciences}, \emph{94}, 785-822.

Light, L., \& Anderson, P. (1993). Designing better keyboards via
simulated annealing. \emph{AI Expert}, \emph{8}(9), 20-27.

Masini, F. (2019). Multi-word expressions and morphology. In
\emph{Oxford Research Encyclopedia of Linguistics}.

McEnery, A., \& Xiao, Z. (2004). The Lancaster Corpus of Mandarin
Chinese: A corpus for monolingual and contrastive language study.
\emph{Religion}, \emph{17}, 3-4.

Noyes, J. (1983). Chord keyboards. \emph{Applied Ergonomics},
\emph{14}(1), 55-59.

Onsorodi, A. H. H., \& Korhan, O. (2020). Application of a genetic
algorithm to the keyboard layout problem {[}Article{]}. \emph{Plos One},
\emph{15}(1), 11, Article e0226611.
\url{https://doi.org/10.1371/journal.pone.0226611}

Pausch, R., Leatherby, J. H., \& Hall, T. (1991). \emph{Voice input vs.
keyboard accelerators: a user study}. Citeseer.

Rosenberg, R., \& Slater, M. (1999). The chording glove: a glove-based
text input device. \emph{IEEE Transactions on Systems, Man, and
Cybernetics, Part C (Applications and Reviews)}, \emph{29}(2), 186-191.

Rosenberg, R. F. (1998). \emph{Computing without mice and keyboards:
text and graphic input devices for mobile computing}. University of
London, University College London (United Kingdom).

Rude, M. (2019). Using the QWERTY keyboard as a chord keyboard: syllabic
typing by multi-key strokes for language learning \& more. \emph{AI and
machine learning in language education}, 168-180.

Seibel, R. (1962). Performance on a five-finger chord keyboard.
\emph{Journal of Applied Psychology}, \emph{46}(3), 165.

Sharoff, S. (2006). Creating general-purpose corpora using automated
search engine queries. \emph{WaCky}, 63-98.

Tarniceriu, A. D., Dillenbourg, P., \& Rimoldi, B. (2013). The effect of
feedback on chord typing. Proceedings of The Seventh International
Conference on Mobile Ubiquitous Computing, Systems, Services and
Technologies,

\emph{Unicode FAQ}. (2023). Unicode Consortium.
\url{https://unicode.org/faq/}

Wu, F.-G., Chen, C.-Y., \& Chen, C.-H. (2007). Improvements of chord
input devices for mobile computer users. Universal Access in
Human-Computer Interaction. Ambient Interaction: 4th International
Conference on Universal Access in Human-Computer Interaction, UAHCI 2007
Held as Part of HCI International 2007 Beijing, China, July 22-27, 2007
Proceedings, Part II 4,

Zheng, Y., Xie, L., Liu, Z., Sun, M., Zhang, Y., \& Ru, L. (2011). Why
press backspace? Understanding user input behaviors in Chinese Pinyin
input method. Proceedings of the 49th annual meeting of the association
for computational linguistics: Human language technologies,

\begin{enumerate}
\def\labelenumi{\Alph{enumi}.}
\tightlist
\item
\end{enumerate}
